\documentclass[12pt,a4paper]{article}

\usepackage[utf8]{inputenc}
\usepackage[T1]{fontenc}
\usepackage[margin=2.5cm]{geometry}
\usepackage{hyperref}
\usepackage{microtype}
\usepackage{parskip}
\usepackage{lmodern}

\hypersetup{colorlinks=true, urlcolor=blue!60!black}

\pagestyle{empty}

\begin{document}

\begin{flushright}
Kundai Farai Sachikonye \\
Auf der Lichtung 19, 82178 Puchheim, Germany \\
\href{https://github.com/fullscreen-triangle}{github.com/fullscreen-triangle} \\
\texttt{kundai.sachikonye@bitspark.com} \\
(+49) 1626386344 \\[6pt]
18 May 2026
\end{flushright}

\vspace{1em}

Prof.\ Jonas Ries \\
Max Perutz Labs \\
University of Vienna / Medical University of Vienna \\
\texttt{jonas.ries@univie.ac.at}

\vspace{1.5em}

\textbf{Re: PhD Position in Super-Resolution Microscopy (Ref.\ 5654) --- late application}

\vspace{1em}

Dear Prof.\ Ries,

I am writing to express interest in the predoctoral position (Ref.\ 5654) in the
Ries lab. I am aware the advertised deadline has passed and that the position
formally requires passage through the Vienna BioCenter International PhD programme
recruitment --- I am reaching out directly in case the position remains open or in
case you have a related opening. If neither applies now, I would welcome the
opportunity to discuss the VBC selection process or future openings.

The reason I am writing despite the timing is that the lab's research programme ---
developing computational methods for single-molecule localization microscopy
and fitting super-resolution data to geometric structural models --- is precisely
where my own independent research has been converging, from a different starting
point.

\textbf{Geometric model fitting from sparse localization data.}
The central computational problem in SMLM is reconstructing the geometry of a
protein complex from a sparse, noisy cloud of single-molecule positions.
\href{https://github.com/fullscreen-triangle/verum}{Verum} addresses the analogous
problem in a different domain: a vehicle as a 20-node oscillatory circuit graph,
where the complete internal state must be reconstructed from observations at 33\,\%
of nodes (5 out of 15 surface-accessible subsystems). The trajectory completion
operator $T = P \circ B \circ K$ is a Banach contraction with Lipschitz constant
$\lambda \in [0.3, 0.7]$, guaranteeing unique fixed-point convergence to the full
system state; validated on real Formula One telemetry data (convergence residual
$7.2 \times 10^{-9}$).
This maps directly onto the LocMoFit problem: given a localization cloud, recover
the underlying geometric model parameters. The circuit graph topology provides
the structural prior; the contraction mapping provides the convergence guarantee.
The same mathematical architecture --- Kirchhoff propagation from boundary
conditions, backward Viterbi inference, thermodynamic feasibility projection ---
applies when the ``nodes'' are protein subunits and the ``observations'' are
fluorophore positions in a MINFLUX or STORM dataset.

\textbf{Simulator-based inference and surrogate models.}
The lab uses simulator-based (likelihood-free) inference for fluorophore
localisation in 3D. This is exactly the application class that my
\href{https://github.com/fullscreen-triangle/purpose}{Purpose} framework and
\href{https://github.com/fullscreen-triangle/homo-veloce}{homo-veloce} surrogate
model system address. The central result --- the \textbf{Backward Training Principle}
--- establishes that a surrogate trained at the hardest single-objective instance
(the \emph{pure-intent regime}) is competent across every less-constrained
sub-regime by feasibility inclusion, not generalisation. For SMLM simulator-based
inference, the pure-intent regime is the highest-SNR, lowest-density acquisition
condition (one fluorophore per diffraction volume, maximum photon budget) ---
the regime from which all realistic acquisition conditions are feasibility
restrictions. Training the surrogate there eliminates the need to sample the
full acquisition parameter space. A Forward Training Collapse Theorem establishes
a provably positive VC failure-gap for the conventional easy-to-hard approach;
a Sample Complexity Theorem gives $(N+1)\times$ savings for $N$-layer constraint
cascades. The homo-veloce system demonstrates the pattern in practice: Gaussian
process and graph neural network physics solvers distilled into lightweight neural
networks (1D CNNs, LSTMs) for real-time inference --- the same pattern the lab's
deep learning localisation pipeline implements.

\textbf{Multi-modal microscopy data integration.}
\href{https://github.com/fullscreen-triangle/helicopter}{Helicopter} is a
multi-modal microscopy analysis framework that combines fluorescence, brightfield,
phase-contrast, and spectral imaging modalities through sequential constraint
satisfaction. Each physical measurement is treated as an independent exclusion
operator rather than as a signal to be averaged: the framework reduces structural
ambiguity from $\sim 10^{60}$ possible molecular configurations to unique solutions
without depending on the resolution limit of any single modality.
For structural cell biology applications combining SMLM with cryo-EM or confocal
reference channels --- the kind of multi-scale structural integration the lab's
work on clathrin-mediated endocytosis and motor protein stepping requires ---
this is the relevant architecture: each modality contributes independent
geometrical constraints, and the intersection of those constraints determines
the structure.

\textbf{First-principles fluorophore spectroscopy.}
The \href{https://honjo-masamune.vercel.app/tools}{Borgia} framework derives
spectroscopic properties of molecules --- electronic transitions, absorption,
emission, and vibrational modes --- from bounded phase space geometry alone,
with zero empirical parameters. The Triple Equivalence Theorem equates oscillatory
dynamics, categorical state enumeration, and partition operations; a GPU fragment
shader running the partition calculation is formally equivalent to a physical
spectrometer. Validated against NIST for nine elements and six molecules with
self-consistency errors below 1.63\,\%.
For rational fluorophore selection in two-colour or three-colour SMLM, the ability
to compute excitation and emission spectra, extinction coefficients, and Stokes
shifts from molecular structure rather than from empirical measurements enables
screening of candidate labels computationally before synthesis or purchase.
The same framework --- Beer-Lambert absorption derived from Kirchhoff circuit laws
(proved in the
\href{https://github.com/fullscreen-triangle/lagrangian}{Lagrangian} platform,
Resonant Stack theorem, $T = 1/e = 0.3679$ for both) --- provides a first-principles
model of light propagation through biological tissue for thick-sample PSF
characterisation and aberration correction in 3D SMLM.

\textbf{Scientific data querying and localization database search.}
A complementary framework --- the PDSVM (Purpose-Driven Shader Virtual Machine)
--- addresses large-scale scientific database search via bounded oscillatory systems.
Every stable physical object is described by three S-entropy coordinates
$(S_k, S_t, S_e) \in [0,1]^3$ (knowledge, temporal, and evolution entropy);
every query is classified into one of three GPU shader primitive classes
(proved exhaustive by the Completeness Theorem);
and retrieval via a ternary trie over Morton-ordered S-entropy coordinates
runs in $O(k)$ time independent of database size $N$ (373$\times$ speedup at
$N=10^5$, $< 0.01\,\mathrm{ms}$ per iteration on a consumer GPU, WebGPU browser-deployable).
The stable slice $T^*$ is the fixed point of the probe operator $\Pi_P(T)$ ---
the set of data states whose S-entropy coordinates are closest to a given purpose
point $P$ on the Fisher-metric manifold.
This maps exactly onto the LocMoFit problem structure:
the localization cloud is the data set; the target geometric model is the purpose
point $P$; and the stable slice is the protein complex geometry satisfying the
forcing conditions (structural constraints derivable without statistical averaging).
For a lab that builds and extends LocMoFit, the trie gives $O(k)$
localization database search independent of cloud size, and the stable slice
formalism provides a first-principles account of which geometric solutions are
structurally necessary rather than merely statistically likely.
10/10 validation experiments all pass (including generative capability:
1000/1000 instances satisfy forcing conditions within $\varepsilon = 0.05$
of the purpose point).

\textbf{Background.}
My MSc in Computational Biology (Universität Tübingen) and doctoral research
in Computational Lipidomics at TUM give me the biological and computational
science foundation. My primary development environment is Python (PyTorch,
scikit-learn, Ray) and Rust for performance-critical components. I have not
operated STORM or MINFLUX hardware --- I am aware that the PhD programme
provides that training --- but the computational methods I have developed
are directly applicable to the data analysis and method development problems
the lab addresses.
I am legally resident in Germany with full work authorisation and can commute
to Vienna or relocate if needed. English is native; German is conversational.

I recognise this is a late contact and appreciate your time in considering it.

\vspace{1.5em}
Yours sincerely,

\vspace{2.5em}
\textbf{Kundai Farai Sachikonye}

\end{document}
